\documentclass[11pt]{article}
% ===========================================================================
%  Shared preamble for the AP Calculus AB exam-preparation worksheets.
%  Input right after \documentclass. Files beginning with "_" are not compiled
%  on their own. Navy and gold, to match the course's lesson pages.
% ===========================================================================
\usepackage[letterpaper,top=2.3cm,bottom=2.2cm,left=1.7cm,right=1.7cm,headheight=28pt,footskip=30pt]{geometry}
\usepackage{amsmath,amssymb}
\usepackage[T1]{fontenc}
\usepackage[utf8]{inputenc}
\usepackage{lmodern}
\usepackage{xcolor}
\usepackage[most]{tcolorbox}
\usepackage{enumitem}
\usepackage{fancyhdr}
\usepackage{array}
\usepackage{tikz}
\usepackage{pgfplots}
\pgfplotsset{compat=1.18}

\definecolor{navy}{HTML}{161E54}
\definecolor{ink}{HTML}{27357E}
\definecolor{gold}{HTML}{B98A00}
\definecolor{goldbg}{HTML}{FFFDF3}
\definecolor{inkbg}{HTML}{F6F7FD}
\definecolor{rule}{HTML}{C9CFF0}
\definecolor{okgreen}{HTML}{1B7A44}
\definecolor{okbg}{HTML}{F1F8F3}

\pagestyle{fancy}
\fancyhf{}
\fancyhead[L]{\footnotesize Name: \underline{\hspace{4cm}}}
\fancyhead[C]{\textbf{\color{navy}Integration Academy}\\[-2pt]{\scriptsize Calculus AB \textperiodcentered{} AP\textsuperscript{\textregistered} Exam Preparation}}
\fancyfoot[L]{\tiny\color{black!55}\parbox[t]{0.78\textwidth}{AP\textsuperscript{\textregistered} is a registered trademark of the College Board, which was not involved in the production of, and does not endorse, this material. \textcopyright{} Integration Academy.}}
\fancyfoot[R]{\scriptsize\thepage}
\renewcommand{\headrulewidth}{1.2pt}
\renewcommand{\headrule}{\hbox to\headwidth{\color{gold}\leaders\hrule height \headrulewidth\hfill}}
\renewcommand{\footrulewidth}{0pt}
\setlength{\parindent}{0pt}
\setlength{\parskip}{4pt}

% Title block: \wstitle{code}{title}{unit line}
\newcommand{\wstitle}[3]{\begin{center}{\small\bfseries\color{gold}WORKSHEET #1}\\[3pt]{\LARGE\bfseries\color{navy}#2}\\[3pt]{\small\color{black!60}#3}\end{center}\vspace{2pt}}
% Section heading
\newcommand{\sech}[1]{\par\vspace{10pt}{\large\bfseries\color{navy}#1}\par\nobreak\vspace{2pt}{\color{rule}\hrule height0.8pt}\vspace{7pt}}
% A level tag printed beside a group of questions
\newcommand{\level}[1]{\par\vspace{6pt}{\footnotesize\bfseries\color{gold}\MakeUppercase{#1}}\par\nobreak\vspace{2pt}}

\newtcolorbox{ideabox}{colframe=ink,colback=inkbg,coltitle=white,colbacktitle=ink,fonttitle=\bfseries\small,title={Key ideas},arc=2pt,boxrule=0.8pt,left=8pt,right=8pt}
\newtcolorbox{example}[1]{colframe=rule,colback=white,coltitle=navy,colbacktitle=inkbg,fonttitle=\bfseries\small,title={#1},arc=2pt,boxrule=0.8pt,borderline west={3pt}{0pt}{ink},left=10pt,right=8pt}
\newtcolorbox{frbox}[1]{colframe=gold,colback=goldbg,coltitle=white,colbacktitle=gold,fonttitle=\bfseries\small,title={#1},arc=2pt,boxrule=0.8pt,breakable,left=8pt,right=8pt}
\newtcolorbox{solbox}{colframe=okgreen,colback=okbg,coltitle=white,colbacktitle=okgreen,fonttitle=\bfseries\small,title={Worked solutions},arc=2pt,boxrule=0.8pt,breakable,left=8pt,right=8pt}
\newtcolorbox{scorebox}[1]{colframe=navy,colback=white,coltitle=white,colbacktitle=navy,fonttitle=\bfseries\small,title={#1},arc=2pt,boxrule=0.8pt,breakable,left=8pt,right=8pt}

\newcommand{\soln}{\par\smallskip\textbf{\color{ink}Solution.}\ }
\newcommand{\result}[1]{\par\smallskip\textbf{#1}}
% blank working space of a given height
\newcommand{\work}[1]{\par\vspace{#1}}
% a point value in the margin of a scoring guide
\newcommand{\pts}[1]{\hfill{\footnotesize\bfseries\color{gold}[#1]}}

% Coordinate grid for graph questions: \begin{apgraph}{xmin}{xmax}{ymin}{ymax} ... \end{apgraph}
\newenvironment{apgraph}[4]{\begin{center}\begin{tikzpicture}\begin{axis}[
  axis lines=middle, xmin=#1, xmax=#2, ymin=#3, ymax=#4,
  grid=both, grid style={black!12}, xtick distance=1, ytick distance=1,
  ticklabel style={font=\scriptsize}, xlabel={$x$}, ylabel={$y$},
  width=10cm, height=7.2cm, samples=80, clip=false,
  every axis plot/.append style={line width=1.1pt, color=ink}]}{\end{axis}\end{tikzpicture}\end{center}}
\newcommand{\closeddot}[2]{\addplot[only marks, mark=*, mark size=2.2pt, color=ink] coordinates {(#1,#2)};}
\newcommand{\opendot}[2]{\addplot[only marks, mark=*, mark size=2.2pt, color=ink, mark options={fill=white}] coordinates {(#1,#2)};}

% Multiple-choice options: four across, or stacked when the options are long.
\newcommand{\opts}[4]{\par\vspace{3pt}\begin{tabular}{@{}p{0.235\linewidth}p{0.235\linewidth}p{0.235\linewidth}p{0.235\linewidth}@{}}(A)\ #1 & (B)\ #2 & (C)\ #3 & (D)\ #4\end{tabular}\par\vspace{9pt}}
\newcommand{\optsv}[4]{\par\vspace{3pt}\begin{tabular}{@{}l@{\ \ }p{0.9\linewidth}@{}}(A) & #1\\(B) & #2\\(C) & #3\\(D) & #4\end{tabular}\par\vspace{9pt}}

\fancyhead[R]{\footnotesize\color{black!60}Worksheet 6.1}
\begin{document}

\wstitle{6.1}{Accumulating Change}{Unit 6: Integration and Accumulation of Change}

\begin{ideabox}
\begin{itemize}[leftmargin=*,itemsep=2pt,topsep=2pt]
  \item A rate multiplied by a time gives an amount. On a graph of a rate, an amount is an \textbf{area}.
  \item The area under a rate graph from $t=a$ to $t=b$ is the \textbf{net change} in the quantity over that interval.
  \item Area below the axis counts as negative: the quantity is decreasing there.
  \item Amount at time $b$ $=$ amount at time $a$ $+$ net change from $a$ to $b$.
\end{itemize}
\end{ideabox}

\sech{Worked examples}
\begin{example}{Example 1 \textperiodcentered{} A constant rate}
A cyclist rides at a constant 12 km/h for 2.5 hours. How far does she travel?\soln The area of a rectangle of height 12 and width 2.5.\result{30 km.}
\end{example}

\begin{example}{Example 2 \textperiodcentered{} A linear rate}
Water flows at $r(t)=4+2t$ litres per minute for $0\le t\le3$. How much water flows?\soln A trapezoid with parallel sides 4 and 10 and width 3: $\dfrac12(4+10)(3)$.\result{21 litres.}
\end{example}

\begin{example}{Example 3 \textperiodcentered{} A rate that changes sign}
A particle has velocity $v(t)=8-4t$ m/s for $0\le t\le3$. Find its net change in position and the total distance it travels.\soln A triangle of area 8 above the axis on $[0,2]$ and one of area 2 below it on $[2,3]$.\result{Net change 6 m; total distance 10 m.}
\end{example}

\newpage
\sech{Practice}
Questions 1 to 7 use this graph. Water flows into a tank at the rate $r(t)$ litres per minute, where $t$ is in minutes.
\begin{apgraph}{-0.5}{9}{-0.5}{7}
\addplot[color=ink, mark=none, line width=1.1pt] coordinates {(0,2) (2,6) (5,6) (8,0)};
\closeddot{0}{2}\closeddot{8}{0}
\node[font=\small, color=ink] at (axis cs:6.9,3.6) {$y=r(t)$};
\end{apgraph}
\level{Core}
\begin{enumerate}[leftmargin=*,itemsep=2pt]
  \item How much water flows into the tank from $t=0$ to $t=2$?\work{1.1cm}
  \item How much water flows into the tank from $t=2$ to $t=5$?\work{1.1cm}
  \item How much water flows into the tank from $t=0$ to $t=8$?\work{1.1cm}
  \item State the units of the area under the graph of $r$.\work{1.1cm}
\end{enumerate}

\level{Exam level}
\begin{enumerate}[leftmargin=*,itemsep=2pt,start=5]
  \item The tank holds 10 litres at $t=0$. How much does it hold at $t=5$?\work{2.1cm}
  \item At what time have exactly 17 litres flowed into the tank?\work{2.1cm}
  \item Find the average rate at which water flows into the tank over $0\le t\le8$.\work{2.1cm}
  \item A particle has velocity $v(t)=10-2t$ m/s for $0\le t\le8$. Find its displacement and the total distance travelled.\work{2.1cm}
\end{enumerate}

\begin{frbox}{Question 9 \textperiodcentered{} Free-response practice}
A factory makes parts at the rate $p(t)$ parts per hour during a 10-hour shift. The graph of $p$ consists of the line segments joining $(0,40)$, $(4,60)$, $(6,60)$ and $(10,20)$.
\begin{enumerate}[label=(\alph*),leftmargin=*,itemsep=2pt]
  \item How many parts are made in the first 4 hours?\work{1.9cm}
  \item How many parts are made during the whole shift?\work{1.9cm}
  \item Interpret $p(6)=60$, and find the time at which the 260th part is completed.\work{1.9cm}
\end{enumerate}
\end{frbox}

\newpage
\begin{solbox}
\begin{enumerate}[leftmargin=*,itemsep=4pt]
  \item A trapezoid: $\dfrac12(2+6)(2)=8$ litres.
  \item A rectangle: $6\times3=18$ litres.
  \item $8+18+\dfrac12(6)(3)=35$ litres.
  \item (litres per minute) $\times$ (minutes) $=$ litres.
  \item $10+8+18=36$ litres.
  \item 8 litres by $t=2$; 9 more at 6 litres per minute takes 1.5 minutes: $t=3.5$.
  \item $\dfrac{35}{8}=4.375$ litres per minute.
  \item Above the axis on $[0,5]$: 25. Below on $[5,8]$: 9. Displacement 16 m; distance 34 m.
  \item \begin{enumerate}[label=(\alph*),leftmargin=*,itemsep=3pt]
    \item $\dfrac12(40+60)(4)=200$ parts.
    \item $200+60(2)+\dfrac12(60+20)(4)=200+120+160=480$ parts.
    \item Six hours into the shift, parts are being made at 60 per hour. 200 parts are done by $t=4$; 60 more at 60 per hour takes one hour: $t=5$.
  \end{enumerate}
\end{enumerate}
\end{solbox}

\end{document}
